\chapter{Identity, Reuse, And Cache}

\section{Bare OID Is Not Enough}

\path{docs/research/experiments/EX-06-identity-tracking/README.md} showed that
the FS root tree logical OID stayed stable while physical address, object XID,
checksum, and block hash changed across mutations. Therefore a cache keyed only
by OID is unsafe.

The relevant synthesis log is
\path{docs/research/RL-04-node-identity-cache-keys-and-oid-reuse.md}.

\section{Subtree Reuse}

\path{docs/research/experiments/EX-07-subtree-reuse/README.md} found zero false
reuse for exact node-identity matches in a detached image-backed lab corpus.
That result keeps subtree summary reuse alive as a future architecture track.

It is not yet a production cache theorem. A future cache entry must include at
least:

\begin{itemize}
  \item source or volume identity
  \item parser version
  \item summary schema version
  \item OMAP domain
  \item OID
  \item object XID
  \item physical address
  \item checksum or content hash
  \item expected type/subtype
\end{itemize}

\section{Persistence And Invalidation}

Persistent cache storage is intentionally deferred. The project must first
replace proof-backend raw walking with native summaries and then validate
incremental output against a fresh full raw parse.

Related logs:

\begin{itemize}
  \item \path{docs/research/RL-05-subtree-reuse-correctness.md}
  \item \path{docs/research/RL-09-cache-persistence-and-invalidation.md}
  \item \path{docs/research/RL-12-performance-model-and-optimization.md}
\end{itemize}
